\subsection{Properties preserved under passage to a direct limit}
\label{subsection:IV.5.13}

\oldpage[IV-2]{131}

Throughout this subsection,
$(A_\alpha,\varphi_{\beta\alpha})$ denotes a direct system of
rings whose index set is filtered and preordered.  We put
\[
  A=\varinjlim_\alpha A_\alpha
\]
and denote by $\varphi_\alpha:A_\alpha\to A$ the canonical
homomorphism.

\begin{proposition}[5.13.1]
\label{IV.5.13.1}
\begin{enumerate}[{\rm(i)}]
  \item For every index $\alpha$, let $\mathfrak a_\alpha$ be an
  ideal of $A_\alpha$ such that
  \[
    \varphi_{\beta\alpha}(\mathfrak a_\alpha)
    \subset\mathfrak a_\beta
    \qquad(\alpha\leq\beta).
  \]
  Then the $\mathfrak a_\alpha$, with the restrictions of the
  $\varphi_{\beta\alpha}$, form a direct system;
  $\mathfrak a=\varinjlim\mathfrak a_\alpha$ is canonically
  identified with an ideal of $A$, and
  \[
    A/\mathfrak a
    \simeq
    \varinjlim_\alpha(A_\alpha/\mathfrak a_\alpha).
  \]
  If, in addition,
  \[
    \mathfrak a_\alpha
    =\varphi_{\beta\alpha}^{-1}(\mathfrak a_\beta)
    \qquad(\alpha\leq\beta),
  \]
  then
  \[
    \mathfrak a_\alpha=\varphi_\alpha^{-1}(\mathfrak a)
    \qquad\text{for every }\alpha.
  \]

  \item Conversely, for every ideal $\mathfrak a$ of $A$, if
  \[
    \mathfrak a_\alpha=\varphi_\alpha^{-1}(\mathfrak a)
    \qquad\text{for every }\alpha,
  \]
  then the $\mathfrak a_\alpha$ form a direct system and
  \[
    \mathfrak a=\varinjlim_\alpha\mathfrak a_\alpha.
  \]
\end{enumerate}
\end{proposition}

\begin{proof}
Assertion {\rm(ii)} is plainly a special case of {\rm(i)}.
In {\rm(i)}, the canonical identification of
$\mathfrak a$ with an ideal of $A$, and that of
$A/\mathfrak a$ with
$\varinjlim(A_\alpha/\mathfrak a_\alpha)$, follow from the
exactness of the direct-limit functor in the category of abelian
groups.

Moreover, $\mathfrak a$ is the union of the increasing family of
images $\varphi_\alpha(\mathfrak a_\alpha)$.  Thus, if
$x_\alpha\in A_\alpha$ and
$\varphi_\alpha(x_\alpha)\in\mathfrak a$, there are
$\beta\geq\alpha$ and $y_\beta\in\mathfrak a_\beta$ such that
\[
  \varphi_\alpha(x_\alpha)=\varphi_\beta(y_\beta).
\]
For some $\gamma\geq\beta$ one then has
\[
  \varphi_{\gamma\alpha}(x_\alpha)
  =\varphi_{\gamma\beta}(y_\beta)\in\mathfrak a_\gamma.
\]
Under the additional inverse-image hypothesis this implies
$x_\alpha\in\mathfrak a_\alpha$, and proves the last assertion of
{\rm(i)}.
\end{proof}

\begin{corollary}[5.13.2]
\label{IV.5.13.2}
Let $\mathfrak N_\alpha$ be the nilradical of $A_\alpha$.  The
nilradical of $A$ is canonically identified with
$\varinjlim\mathfrak N_\alpha$, and
\[
  A_{\mathrm{red}}
  \simeq
  \varinjlim_\alpha(A_\alpha)_{\mathrm{red}}.
\]
In particular, if every $A_\alpha$ is reduced, then $A$ is
reduced.
\end{corollary}

\oldpage[IV-2]{132}

\begin{proof}
It is clear that
\[
  \varphi_{\beta\alpha}(\mathfrak N_\alpha)
  \subset\mathfrak N_\beta
  \qquad(\alpha\leq\beta),
\]
and that $\varinjlim\mathfrak N_\alpha$ is a nil ideal of $A$,
hence is contained in the nilradical $\mathfrak N$ of $A$.
Conversely, if $x\in\mathfrak N$, then $x^h=0$ for some integer
$h$.  Write $x=\varphi_\alpha(x_\alpha)$ with
$x_\alpha\in A_\alpha$.  The relation
$\varphi_\alpha(x_\alpha^h)=0$ implies that for some
$\beta\geq\alpha$,
\[
  \varphi_{\beta\alpha}(x_\alpha^h)=0.
\]
Putting $x_\beta=\varphi_{\beta\alpha}(x_\alpha)$, we obtain
$x_\beta^h=0$, hence $x_\beta\in\mathfrak N_\beta$, while
$x=\varphi_\beta(x_\beta)$.  This proves the assertion, in view
of \sref{IV.5.13.1}.
\end{proof}

\begin{proposition}[5.13.3]
\label{IV.5.13.3}
\begin{enumerate}[{\rm(i)}]
  \item For every index $\alpha$, let $\mathfrak p_\alpha$ be a
  prime ideal of $A_\alpha$ such that
  \[
    \varphi_{\beta\alpha}(\mathfrak p_\alpha)
    \subset\mathfrak p_\beta
    \qquad(\alpha\leq\beta).
  \]
  Then $\mathfrak p=\varinjlim\mathfrak p_\alpha$ is a prime
  ideal of $A$.  In particular, if all the $A_\alpha$ are
  domains, then $A$ is a domain.

  \item If, in addition,
  \[
    \mathfrak p_\alpha
    =\varphi_{\beta\alpha}^{-1}(\mathfrak p_\beta)
    \qquad(\alpha\leq\beta),
  \]
  then there is a canonical identification
  \[
    A_{\mathfrak p}
    \simeq
    \varinjlim_\alpha(A_\alpha)_{\mathfrak p_\alpha}.
  \]

  \item Suppose that $A$ is local, with maximal ideal
  $\mathfrak p$, and that the homomorphisms
  $\varphi_\alpha:A_\alpha\to A$ are injective.  If
  $\mathfrak p_\alpha=\varphi_\alpha^{-1}(\mathfrak p)$, then
  \[
    A
    \simeq
    \varinjlim_\alpha(A_\alpha)_{\mathfrak p_\alpha},
  \]
  and every canonical homomorphism
  $(A_\alpha)_{\mathfrak p_\alpha}\to A$ is injective.
\end{enumerate}
\end{proposition}

\begin{proof}
The two assertions in {\rm(i)} are equivalent, since
\[
  A/\mathfrak p
  \simeq
  \varinjlim_\alpha(A_\alpha/\mathfrak p_\alpha).
\]
Suppose therefore that all the $A_\alpha$ are domains, and let
$x,y\in A$ satisfy $xy=0$.  For some $\alpha$ one may write
\[
  x=\varphi_\alpha(x_\alpha),
  \qquad
  y=\varphi_\alpha(y_\alpha),
\]
with $x_\alpha,y_\alpha\in A_\alpha$ and
$\varphi_\alpha(x_\alpha y_\alpha)=0$.  Hence there is
$\beta\geq\alpha$ such that, on putting
\[
  x_\beta=\varphi_{\beta\alpha}(x_\alpha),
  \qquad
  y_\beta=\varphi_{\beta\alpha}(y_\alpha),
\]
one has $x_\beta y_\beta=0$.  Since $A_\beta$ is a domain,
$x_\beta=0$ or $y_\beta=0$, and therefore $x=0$ or $y=0$.

For {\rm(ii)}, the hypothesis says that the multiplicative
subsets $A_\alpha-\mathfrak p_\alpha$ form a direct system whose
direct limit is $A-\mathfrak p$.  The local rings
$(A_\alpha)_{\mathfrak p_\alpha}$ therefore form a direct system
with local transition maps, and the canonical homomorphisms to
$A_{\mathfrak p}$ induce a surjection
\[
  \psi:
  \varinjlim_\alpha(A_\alpha)_{\mathfrak p_\alpha}
  \longrightarrow A_{\mathfrak p}.
\]
To prove injectivity, let $x_\alpha\in A_\alpha$ and
$s_\alpha\in A_\alpha-\mathfrak p_\alpha$, and suppose that
\[
  \frac{\varphi_\alpha(x_\alpha)}
       {\varphi_\alpha(s_\alpha)}
  =0
  \quad\text{in }A_{\mathfrak p}.
\]
There is $s'\in A-\mathfrak p$ such that
$s'\varphi_\alpha(x_\alpha)=0$.  After replacing $\alpha$ by a
larger index, we may write
$s'=\varphi_\alpha(s'_\alpha)$ with
$s'_\alpha\in A_\alpha-\mathfrak p_\alpha$.  Then
$\varphi_\alpha(s'_\alpha x_\alpha)=0$, so for some
$\beta\geq\alpha$,
\[
  \varphi_{\beta\alpha}(s'_\alpha x_\alpha)=0
  \quad\text{in }A_\beta.
\]
It follows that the image of $x_\alpha/s_\alpha$ in the direct
limit is zero.

Finally, the first assertion of {\rm(iii)} follows from {\rm(ii)}
and $A_{\mathfrak p}=A$.  Every element of
$A_\alpha-\mathfrak p_\alpha$ maps to a unit of $A$, and the
injectivity of
$(A_\alpha)_{\mathfrak p_\alpha}\to A$ follows from that of
$\varphi_\alpha$.
\end{proof}

\begin{corollary}[5.13.4]
\label{IV.5.13.4}
Suppose that all the $A_\alpha$ are domains and that all the
transition homomorphisms $\varphi_{\beta\alpha}$ are injective.
If $A'_\alpha$ denotes the integral closure of $A_\alpha$, then
\[
  A'=\varinjlim_\alpha A'_\alpha
\]
is the integral closure of $A$.  In particular, if every
$A_\alpha$ is integrally closed, then $A$ is integrally closed.

If the $A_\alpha$ are unibranch local rings
(respectively, geometrically unibranch local rings)
\sref[0]{0.23.2.1}, and the $\varphi_{\beta\alpha}$ are injective
local homomorphisms, then $A$ is unibranch
(respectively, geometrically unibranch).
\end{corollary}

\begin{proof}
Let $K_\alpha$ be the fraction field of $A_\alpha$.  By
\sref{IV.5.13.3}[ii], applied to $\mathfrak p_\alpha=(0)$,
\[
  K=\varinjlim_\alpha K_\alpha
\]
is the fraction field of $A$.  Thus
$A'=\varinjlim A'_\alpha$ is identified with

\oldpage[IV-2]{133}

a subring of $K$ which is plainly integral over $A$.

Conversely, let $z\in K$ satisfy an equation of integral
dependence
\[
  z^n+\sum_{i=1}^{n}c_i z^{\,n-i}=0,
  \qquad c_i\in A.
\]
There are an index $\alpha$, an element $z_\alpha\in K_\alpha$,
and elements $c_{i\alpha}\in A_\alpha$ whose images are
$z$ and the $c_i$, respectively, and such that the same equation
holds in $K_\alpha$.  Hence $z_\alpha\in A'_\alpha$, so
$z\in A'$.  This also proves at once that a direct limit of the
integrally closed rings occurring here is integrally closed.

Suppose now that the $A_\alpha$ are local and that the transition
maps are local.  Then $A$ is local and, if $k_\alpha$ denotes the
residue field of $A_\alpha$,
\[
  k=\varinjlim_\alpha k_\alpha
\]
is the residue field of $A$ \sref[0]{0.10.3.1.3}.  If every
$A'_\alpha$ is local, then $A'$ is local, and hence $A$ is
unibranch.  If, moreover, the residue field $k'_\alpha$ of
$A'_\alpha$ is a radicial extension of $k_\alpha$ for every
$\alpha$, then
\[
  k'=\varinjlim_\alpha k'_\alpha,
\]
the residue field of $A'$, is a radicial extension of $k$.
Thus $A$ is geometrically unibranch.
\end{proof}
